\chapter{绪论}\label{chap:introduction}

本章首先介绍面向网络流量的重复音视频文件实时检测系统的研究背景、意义和系统要求，并对国内外研究现状进行详细分析。最后简要介绍本文的主要组织架构。

\section{研究背景及意义}

随着互联网的发展，网络已经渗入到我们生活的方方面面，越来越多的人使用网络观看和下载音视频节目。根据中国互联网信息中心(CNNIC)发布的第42次《中国互联网络发展状况统计报告》\cite{CNNIC2018}显示。截至2018年6月，如表~\ref{tab:Numberofuser}~所示，我国的网民规模达8.02亿人，其中网络视频用户约为6.09亿人，较半年前增长了5.7\%，网络音乐用户约为5.55亿人，较半年前增长了1.2\%。

\begin{table}[!htbp]
    \bicaption{网民人数以及音视频用户人数}{Number of netizens and number of multimedia netizens}
    \label{tab:Numberofuser}
    \centering
    \footnotesize% fontsize
    \setlength{\tabcolsep}{4pt}% column separation
    \renewcommand{\arraystretch}{1.2}%row space 
    \begin{tabular}{lccc}
        \hline
        &2017.12用户规模(万)&2018.6用户规模(万)&半年增长率\\
        %\cline{2-9}% partial hline from column i to column j
        \hline
        网民总数&80166&77198&3.8\%\\
        网络视频用户&57892&60906&5.2\%\\
        网络音乐用户&54809&55482&1.2\%\\
        \hline
    \end{tabular}
\end{table}

其中网络直播和短视频用户飞速增长，2017年起，尤其在2018年春节期间，短视频应用迅速下沉至三、四线城市，用户规模持续增长。截至2018年6月，综合各个热门短视频应用的用户规模达5.94亿人，网络直播的用户规模也达到4.25亿人，较去年年末增涨294万人。

随着通过网络观看音视频节目的人数不断增加，音视频流量也早已成为网络流量中占比最高的流量，根据Sandvine发布的2018年全球互联网观察\cite{PhenomenaReport}显示，如表~\ref{tab:Differentapplication}~所示，视频流量已经占据全球流量的57.69\%，较去年增长了22.43\%。

\begin{table}[!htbp]
    \bicaption{不同应用网络流量占比}{Global application category traffic share}
    \label{tab:Differentapplication}
    \centering
    \footnotesize% fontsize
    \setlength{\tabcolsep}{4pt}% column separation
    \renewcommand{\arraystretch}{1.2}%row space 
    \begin{tabular}{lcc}
        \hline
        应用名&流量占比&全年增长率\\
        %\cline{2-9}% partial hline from column i to column j
        \hline
        视频流&57.64\%&22.43\%\\
        网页&17.01\%&20.98\%\\
        游戏&7.78\%&2.68\%\\
        社交媒体&5.10\%&3.73\%\\
        电子商务&4.61\%&1.90\%\\
        文件分享&2.84\%&22.05\%\\
        即时通讯&1.72\%&8.12\%\\
        安全&1.41\%&7.48\%\\
        \hline
    \end{tabular}
\end{table}

美国下载流量前十位的应用中，有四种应用提供音视频服务，如表~\ref{tab:Differentdownload}~所示，分别为Netflix，YouTube，HTTP media stream，Amazon Prime。

\begin{table}[!htbp]
    \bicaption{美国下载流量排行}{US application traffic share}
    \label{tab:Differentdownload}
    \centering
    \footnotesize% fontsize
    \setlength{\tabcolsep}{4pt}% column separation
    \renewcommand{\arraystretch}{1.2}%row space 
    \begin{tabular}{lcc}
        \hline
        应用名&流量占比&全年增长率\\
        %\cline{2-9}% partial hline from column i to column j
        \hline
        Netflix&14.97\%&2.92\%\\
        HTTP media stream&13.07\%&4.84\%\\
        YouTube&11.35\%&3.03\%\\
        Raw MPEG-TS&4.39\%&4.11\%\\
        HTTP (TLS)&4.06\%	&2.06\%\\
        QUIC&3.87\%&1.43\%\\
        Amazon Prime&3.69\%&0.87\%\\
        HTTP download&3.69\%&1.45\%\\
        HTTP&3.22\%&4.80\%\\
        PlayStation download&2.67\%&0.45\%\\
        \hline
    \end{tabular}
\end{table}

根据\cite{CiscoReport}，到2022年全球音视频流量将占到全球网络流量的82\%,高于2017年的75\%。从2017年到2022年，互联网视频流量将增长四倍，复合年增长率约为33\%。

根据\cite{Yu}发表的论文，发现大部分的视频用户会重复下载观看同一部电影，且YouTube流量中约30\%的重复，如图~\ref{fig:youtube}~所示。可见，随着音视频流量的不断增多，大量重复的音视频流量通过网络进行传播。
\begin{figure}[!htbp]
    \centering
    \includegraphics[width=0.80\textwidth]{youtube}
    \bicaption{PACK网络冗余检测结果}{PACK redundancy elimination ratio}
    \label{fig:youtube}
\end{figure}

代理服务器，入侵检测系统等网络中间件需要通过处理分析流量，以达到版权保护、色情检测等目的。然而对这些重复流量的处理消耗了这些网络中间件大量带宽，存储和处理资源，同时也降低了整个网络性能。如果这些网络中间件能在音视频流量建立连接时，对这些重复的音视频文件进行检测，就可以提高网络性能并节约服务器资源。

综上所述，在网络流中识别重复的音视频文件具有重要意义。针对网络流量中音视频文件的重复传输问题，本文提出一个准确性高的、实时的重复音视频文件实时检测系统。

\section{系统要求}
面向网络流量的重复音视频文件实时检测系统，其有两个要求：

1.检测的实时性：为了在网络流量中检测重复的音视频文件，从而减少代理服务器等网络中间件的消耗，要尽可能在传输结束前完成检测。选取的特征应主要为音视频流量头部的信息，以保证检测的实时性。

2.算法的准确性：由于代理服务器入侵检测系统等网络中间件需要对网络流量进行分析从而达到版权保护和色情检测的目的。因此对于音视频文件的重复性检测，其对准确性的要求较高。

\section{国内外研究现状}
\subsection{研究现状概述}
在网络流量中对重复的音视频文件实时检测主要可以分为两种方法:

第一种只分析网络流量，不进行音视频文件的编解码，即网络冗余发现系统，现大量应用于云计算，企业内部网络等，其优点是大大减少了网络的建设成本，缺点是无法检测相似的网络流量。 如：\cite{Spring2000}提出一种无需对网络协议进行解析的网络冗余发现系统，该论文首次提出一种端到端的模型，并使用拉宾指纹对传输的流量分块。

还有一些网络冗余发现系统带有一定的预测功能。\cite{Zohar2014}提出一种具有预测性的网络冗余发现算法,该方法与前人提出的基于端到端网络冗余发现系统的不同点在于，该方法不要求服务器持续维护客户端状态（紧密的缓存同步问题），大大地减少了服务器的开销。

网络冗余发现系统被用于不同的网络细分领域，与使用的场景密切结合，用于减少该场景下资源的消耗。\cite{Lei2013Cooperative}提出了一种用于云计算中的网络冗余发现，该系统即能够发现短冗余又能发现长冗余。\cite{Aggarwal2010EndRE}提出的网络冗余发现系统被用于广域网的优化。\cite{Sanadhya2012}将现有的网络冗余发现技术用于手机中。\cite{Halepovic2012}在数据链路层解决无线局域网中的网络冗余问题，在该论文中也首次提出在无线网络中发现网络冗余的四个挑战。

第二种是对音视频文件编解码，即大规模相似音视频检索（ NDVR ）系统，现大量应用于版权保护，特定视频检测等领域。其优点是检测的准确率较高，缺点是消耗大量的资源，无法实时地检测视频的相似性。如\cite{Hao2016Stochastic}为了提高检索相似音视频的准确性和速度，提出了一种新的随机多视图哈希算法，以促进大规模相似音视频检索（ NDVR ）系统的构建。最后使用两个公开的相似音视频数据集对该系统进行检验，与各种最先进的NDVR系统进行比较有着较明显的优势。

下面我将按照网络的分层结构，详细介绍网络冗余发现系统。
\subsection{数据链路层的网络冗余发现系统}
2010年SIGCOMM上，\cite{Lumezanu2010The}等人发表的一篇论文提出了一种针对蜂窝数据网络的冗余发现和消除方法。在该篇论文中，作者在蜂窝络中部署网络冗余发现系统，可以大大减少运营商的网络带宽成本。由于蜂窝链路具有较高的丢包率，给整个网络冗余发现系统带了巨大的挑战。该论文作者提出了一种轻量级机制，来应对此问题。作者将该系统部署在两个北美和一个欧洲运营商处，实验表明该网络冗余发现系统可以将大多数移动用户的带宽消耗降低至少约10％，其中部分带宽消耗减少高达50％，取得了良好的效果。

\cite{Halepovic2012}在2012年ICC上发表的一篇论文提出了一种针对无线局域网链路层的网络冗余发现系统。在论文中，作者证明了在无线局域网中部署网络冗余发现系统是可行的，并将其和使用在以太网中的网络冗余发现系统进行了比较。作者提出了无线局域网网络冗余发现系统所面临的两个挑战：一是链路的不稳定性，在数据链路层会出现大量的帧丢失现象。二是无线局域网的数据包中包含大量管理和控制帧，导致IP数据的比例下降。作者通过实验证明，虽然网络消冗系统在无线局域网中的效果并没有在以太网中那么明显，但依旧可以节约23\%的带宽消耗。主要的原因是数据链路层需要对丢失的帧进行再重传，产生了大量的额外流量。

\subsection{网络层的网络冗余发现系统}

为了避免对大量的帧丢失进行处理，大量的网络冗余发现系统都工作在网络层和应用层。\cite{Spring2000}提出了一种识别重复数据传输的技术，并用它来分析网络流量中的冗余。整个系统分成发送端、接收端和存储数据块及其指纹的数据集，指纹数据集需要在发送端和接收端保持同步。该系统是一个典型的基于端到端模型的网络冗余发现系统,系统所使用的模型和算法思想也影响了许多后人的工作。发送端的流程如下：第一步：发送端对网络流量基于内容进行分块，使用了滑动窗口技术，用拉宾指纹算法来确定网络流量分块的边界。使用拉宾指纹的目的是为了避免不变的分块对数据的更改十分敏感。其原理如下图~\ref{fig:Content-based}~所示。第二步：每次计算给定滑动窗口W内的数据内容的指纹FP（如W），并基于最大给定的大整数D取模后，与给定的余数进行了对比，若相等（如W2）则窗口的右端为数据划分边界，否则（如W1）将窗口继续往右滑动一个字节，依此循环地进行计算和对比，直到达到文件末尾。基于内容分块可以将更改对数据划分的影响控制在更新位置附近的少数几个分块内，并保持其他分块不变。第三步：发送端将分块后的数据在数据集中进行查询，若在数据集中查询到该分块，则将该分块替换成其在数据集中对应的指纹。若在数据集中查询不到，则使用新的指纹对需要传输的数据块进行标记，并存储在数据集中。接着将新指纹和需要传输的数据块一起发给接收端，以保证接收端的数据集的更新。接收端的流程如下：接收端收到发送端的网络流量后，需要对网络流量进行还原，若是新指纹，则记录在数据集中。若是旧指纹，则对指纹进行还原，填充相应的数据块。
\begin{figure}[!htbp]
    \centering
    %trim option's parameter order: left bottom right top
    \includegraphics[width=0.80\textwidth]{Content-based}
    \bicaption{基于内容分块原理}{Content-Defined Chunking approach}
    \label{fig:Content-based}
\end{figure}

\cite{Yu}提出了PACK压缩算法，实现了一个专门为云计算设计的网络冗余发现系统。整个系统也采用了典型的基于端到端的模型，但本文提出了一种新的网络冗余发现算法PACK，不同于以往的算法，PACK有一定的预测性质。PACK也减少了发送端的开销，无需发送端和接收端紧密缓存两个数据集。简单介绍一下PACK算法。PACK接收端算法步骤如下，当接收端接收到发送端发送的数据分块，会在数据集中查询。若在数据集中查询到该数据分块，则返回数据集中该数据分块的下一块数据指纹，即PRED，并等待发送端回复PRED-ACK。若发送端回复PRED-ACK，则认为下一块数据已经接收，以此类推。若在数据集中查询不到，告知发送端继续发送。PACK发送端算法如下，发送端将需要发送的信息用弱哈希进行分块，并生成指纹，在发送过程中会不断接收PRED，将其中包含的指纹与将要发送的分块数据的指纹进行对比，若相同，则返回PRED-ACK。若不相同，则返回数据块。PACK的启发是，文件的重复是可预测的，文件的部分内容和整个文件是有相关性的，可用文件的部分内容标识整个文件。

\subsection{应用层的网络冗余发现系统}

在2012年的NSDI上，\cite{Aggarwal}发表一篇文章，提出一种不同于中间件的网络冗余发现系统，可适用于加密流量和最后一跳是移动设备的应用场景，使用了一种新的优于拉宾指纹的算法。该算法无论在缓存开销方面还是在CPU占用方面都优于以往的算法，缺点是无法发现较短的冗余。该系统需要在用户的电子设备额外安装一个客户端，客户端加解密工作是由原本的TLS协议完成，因此该系统只是适用于企业内部网络，以减少企业在网络建设上的消耗。该篇论文中提到了基于中间件的网络冗余发现系统的缺点，主要包含以下两点：1)随着HTTPS的普及，越来越多的网站使用全站加密，基于中间件的网络冗余发现系统，不能很好地应对加密流量。2)随着移动互联网的兴起，越来越多的人使用手机，基于中间件的网络冗余发现系统无法提高移动设备中最后一跳链路的性能。

从2012年起，对移动设备的网络冗余发现也越来越得到了人们的关注。\cite{Feng2013How}发表一篇文章提出了一种组合不同的网络冗余发现和消除技术的方法，如：增量编码，网络缓存等。在移动网络中减少了约30\%的流量。本文的作者调查了移动设备上的HTTP缓存使用情况，发现移动设备在Web缓存，协议规范等方面与电脑设备之间存在巨大差距。完全支持HTTP缓存协议并且严格遵循该协议的智能手机应用程序和移动浏览器，可以减少约17\%左右的移动流量。

\subsection{国内外研究现状小结}
国内外研究现状的小结，见表~\ref{tab:Summaryofresearch}~所示，这些研究成果要么无法对重复的网络流量进行实时检测，要么都是基于端到端系统，无法在网络中间件中使用。基于音视频内容的NDVR系统消耗大量的资源，且实时性较差，并不适合在网络传输环境中使用。但这些研究成果却有许多可借鉴之处：一是系统应充分考虑网络传输中丢包，乱序等特性对检测结果的影响。二是系统应充分考虑不同网络设备的适用。三是系统应充分考虑内存消耗的问题，可以在内存集约的情况下工作。四是文件的部分内容和整个文件是有相关性的，可用文件的部分内容标识整个文件。
\begin{table}[!htbp]
    \bicaption{国内外研究现状小结}{Summary of research status}
    \label{tab:Summaryofresearch}
    \centering
    \footnotesize% fontsize
    \setlength{\tabcolsep}{4pt}% column separation
    \renewcommand{\arraystretch}{1.2}%row space 
    \begin{tabularx}{\textwidth}{lccX}
        \hline
        研究方法&代表工作&会议&创新点\\
        \hline
        \multirow{2}*{链路层}&RTE in WLAN&ICC'12&解决了数据帧丢失和IP数据比例下降问题。\\
        ~ & Effect of Packet loss & SIGCOMM'10 & 首次将网络冗余发现系统应用于蜂窝数据网络。 \\
        \multirow{2}*{网络层}&Protocol-independent RE Algorithms&SIGCOMM'00&首次提出基于端到端的网络冗余发现系统，且可应用于任何非加密流量。\\
        ~ & end to end TRE & ICNP'12 & 可适用于不同网络设备，专门为云计算设计的网络冗余检测和消除算法。 \\
        \multirow{3}*{应用层}&EndRE&NSDI'12&适用于加密流量和最后一跳是移动设备的应用场景，同时使用了一种新的优于拉宾指纹的算法。\\
        ~ &Combination of different techniques&PAM'13&组合不同的网络冗余发现和消除技术的方法，如：增量编码，网络缓存等。在移动网络中减少了约30\%的流量。\\
        ~ &PACK&TON'14&提出了一种具有预测性的网络冗余检测系统，大大减少了服务器端的压力。\\
        基于内容&NDVR&TMM'17&提出了一种新的随机多视图哈希算法，以促进大规模相似音视频检索（ NDVR ）系统的构建。\\
        \hline
    \end{tabularx}
\end{table}
% 最后作者采用了现有的四种冗余发现方法，分别采用了HTTP缓存，文件压缩，增量编码和数据包流压缩(MODP),最终减少了约30\%的手机流量。
% 其结果如图\ref{fig:Mobiledevice}所示。

% \begin{figure}[!htbp]
%     \centering
%     \begin{subfigure}[a]{0.90\textwidth}
%       \includegraphics[width=\textwidth]{Mobiledevice1}
%       \caption{}
%       \label{fig:Mobiledevice1}
%     \end{subfigure}%
%     \\~% add desired spacing
%     \begin{subfigure}[b]{0.80\textwidth}
%       \includegraphics[width=\textwidth]{Mobiledevice2}
%       \caption{}
%       \label{fig:Mobiledevice2}
%     \end{subfigure}
%     \bicaption{移动设备流量减少结果}{Mobilephone traffic reduce results}
%     \label{fig:Mobiledevice}
% \end{figure}
\section{论文的主要内容与组织结构}
\subsection{论文的主要内容}
本文的研究目标：针对网络中间件消耗大量资源处理重复音视频文件的问题，设计并实现面向网络流量的重复音视频实时检测系统。

本文实现的系统包含两个关键技术：1）预测性文件标识生成方法。该技术在流量建立连接时，生成一个具有预测性的文件标识。主要步骤如下：选取流量头部的一些特征，通过MD5等哈希算法，生成文件标识。2）基于流式模糊哈希的重复音视频检测方法。主要步骤如下：在数据集中查询文件标识，同时计算流量的流式模糊哈希值，用于验证标识的准确性，指导标识的生成。

实验证明，该系统可以去除约60\%的重复文件，不影响内容分析模块特定音视频文件检测的准确率。
\subsection{组织结构}
论文内容共分为六章，每章的安排如下：

第1章 绪论。本章简要介绍了论文的背景和研究意义，对国内外的研究现状进行全面详细分析，最后简要介绍本文的研究内容和组织结构。

第2章 前人的工作基础。本章介绍了流式模糊哈希和基于流式模糊哈希的相似性查找系统的原理，并分析各自的使用原因，和优点。

第3章 预测性文件标识生成方法。本章介绍了预测性文件标识生成方法，通过特征筛选生成了6个候选的文件标识。

第4章 基于流式模糊哈希的重复音视频检测方法。本章介绍了该方法的步骤，通过实验证明了该方法的可行性，实验的结果指导了文件标识的生成。

第5章 系统设计与实现。本章介绍了重复音视频实时检测系统的设计与实现，并在某公司网关对整个系统进行了测试。

第6章 总结与展望。本章对本文的工作进行全面总结，并对未来的工作进行了展望。

\section{小结}
本章首先介绍了本文的研究背景和意义，大量重复的音视频文件通过网络进行传输，对这些重复的文件进行检测可以减少网络中间件的消耗，提高整个网络性能。接着提出了重复音视频实时检测系统的两个要求：一是检测实时性。二是算法准确性。然后介绍了国内外的研究现状，并按网络分层的顺序对网络冗余发现系统进行了详细介绍。最后简要介绍了本文的研究内容和组织结构。